\documentclass[10pt,a4paper]{amsart}
\usepackage{amsmath,amssymb,amsthm,graphicx}
\usepackage{verbatim}
\newenvironment{reference}{\comment}{\endcomment}
%\newenvironment{reference}{}{}
\newenvironment{slogan}{\comment}{\endcomment}
\newenvironment{history}{\comment}{\endcomment}

% For commutative diagrams we use Xy-pic
\usepackage[all]{xy}

% We use 2cell for 2-commutative diagrams.
\xyoption{2cell}
% Theorem environments.
%
\theoremstyle{plain}
\newtheorem{theorem}[subsection]{Theorem}
\newtheorem{proposition}[subsection]{Proposition}
\newtheorem{lemma}[subsection]{Lemma}

\theoremstyle{definition}
\newtheorem{definition}[subsection]{Definition}
\newtheorem{example}[subsection]{Example}
\newtheorem{exercise}[subsection]{Exercise}
\newtheorem{situation}[subsection]{Situation}

\theoremstyle{remark}
\newtheorem{remark}[subsection]{Remark}
\newtheorem{remarks}[subsection]{Remarks}

\numberwithin{equation}{subsection}

% Macros
%
\def\lim{\mathop{\mathrm{lim}}\nolimits}
\def\colim{\mathop{\mathrm{colim}}\nolimits}
\def\Spec{\mathop{\mathrm{Spec}}}
\def\Hom{\mathop{\mathrm{Hom}}\nolimits}
\def\Ext{\mathop{\mathrm{Ext}}\nolimits}
\def\SheafHom{\mathop{\mathcal{H}\!\mathit{om}}\nolimits}
\def\SheafExt{\mathop{\mathcal{E}\!\mathit{xt}}\nolimits}
\def\Sch{\mathit{Sch}}
\def\Mor{\mathop{\mathrm{Mor}}\nolimits}
\def\Ob{\mathop{\mathrm{Ob}}\nolimits}
\def\Sh{\mathop{\mathit{Sh}}\nolimits}
\def\NL{\mathop{N\!L}\nolimits}
\def\CH{\mathop{\mathrm{CH}}\nolimits}
\def\proetale{{pro\text{-}\acute{e}tale}}
\def\etale{{\acute{e}tale}}
\def\QCoh{\mathit{QCoh}}
\def\Ker{\mathop{\mathrm{Ker}}}
\def\Im{\mathop{\mathrm{Im}}}
\def\Coker{\mathop{\mathrm{Coker}}}
\def\Coim{\mathop{\mathrm{Coim}}}

% Boxtimes
%
\DeclareMathSymbol{\boxtimes}{\mathbin}{AMSa}{"02}

%
% Macros for moduli stacks/spaces
%
\def\QCohstack{\mathcal{QC}\!\mathit{oh}}
\def\Cohstack{\mathcal{C}\!\mathit{oh}}
\def\Spacesstack{\mathcal{S}\!\mathit{paces}}
\def\Quotfunctor{\mathrm{Quot}}
\def\Hilbfunctor{\mathrm{Hilb}}
\def\Curvesstack{\mathcal{C}\!\mathit{urves}}
\def\Polarizedstack{\mathcal{P}\!\mathit{olarized}}
\def\Complexesstack{\mathcal{C}\!\mathit{omplexes}}
% \Pic is the operator that assigns to X its picard group, usage \Pic(X)
% \Picardstack_{X/B} denotes the Picard stack of X over B
% \Picardfunctor_{X/B} denotes the Picard functor of X over B
\def\Pic{\mathop{\mathrm{Pic}}\nolimits}
\def\Picardstack{\mathcal{P}\!\mathit{ic}}
\def\Picardfunctor{\mathrm{Pic}}
\def\Deformationcategory{\mathcal{D}\!\mathit{ef}}
\usepackage{cleveref}
\crefname{originalreference}{original reference}{original references}
\Crefname{originalreference}{Original reference}{Original references}
\makeatletter
\newlabel{lemma-bound-geometric-genus}{{0CE1}{source}{}{sourcepointer-1}{}}
\newlabel{lemma-bound-geometric-genus@cref}{{[originalreference][0][]0CE1}{[1][0][]source}{}{}{}}
\makeatother
\crefname{originallemma}{original lemma}{original lemmas}
\Crefname{originallemma}{Original lemma}{Original lemmas}
\makeatletter
\newlabel{morphisms-lemma-normalization-birational}{{0BXC}{source}{}{sourcepointer-2}{}}
\newlabel{morphisms-lemma-normalization-birational@cref}{{[originallemma][0][]0BXC}{[1][0][]source}{}{}{}}
\makeatother
\crefname{originallemma}{original lemma}{original lemmas}
\Crefname{originallemma}{Original lemma}{Original lemmas}
\makeatletter
\newlabel{varieties-lemma-finite-in-codim-1}{{0AB7}{source}{}{sourcepointer-3}{}}
\newlabel{varieties-lemma-finite-in-codim-1@cref}{{[originallemma][0][]0AB7}{[1][0][]source}{}{}{}}
\makeatother
\crefname{originallemma}{original lemma}{original lemmas}
\Crefname{originallemma}{Original lemma}{Original lemmas}
\makeatletter
\newlabel{varieties-lemma-modification-normal-iso-over-codimension-1}{{0BFP}{source}{}{sourcepointer-4}{}}
\newlabel{varieties-lemma-modification-normal-iso-over-codimension-1@cref}{{[originallemma][0][]0BFP}{[1][0][]source}{}{}{}}
\makeatother
\title{Invariance of the geometric genus under reduction, birational modification and normalization}
\author{The Stacks Project Authors}
\date{stacks-project:curves.tex:local-snapshot-20261012; GNU FDL 1.2 or later}
\begin{document}
\maketitle
\noindent\textit{Editorial note.} All delivered mathematical text is the author's own, copied from the pinned Stacks Project chapter file curves.tex; no statement or proof step was rewritten, reordered or removed.\par\medskip
\noindent\textit{Editorial note.} The counted proof is the author's complete argument: the reduction invariance, the birational component comparison with function fields and models, and the normalization composition are all retained verbatim.\par\medskip
\noindent\textit{Editorial note.} The label lemma-geometric-genus-normalization names the counted statement, and the proof environment immediately following it in the source is delivered verbatim as its proof.\par\medskip
\noindent\textit{Editorial note.} The counted statement is carried by the author's own lemma environment in the chapter source.\par\medskip
\noindent\textit{Allowed standard prerequisites (editorial).}\begin{itemize}
\item The geometric genus via the cohomological bound. Same chapter, Lemma lemma-bound-geometric-genus. Underlies part (1), the equality of geometric genera of X and X\_red.
\item Birational finite morphisms. Varieties, Lemmas varieties-lemma-finite-in-codim-1 and varieties-lemma-modification-normal-iso-over-codimension-1; Morphisms, Lemma morphisms-lemma-normalization-birational. Makes a proper birational morphism finite and iso over a dense open, and records that normalization is birational.
\end{itemize}\medskip
\ifdefined\hypertarget\hypertarget{sourcepointer-1}{}\fi
\noindent\textit{Source reference.} Original reference 0CE1: sources/stacks-curves/pinned-original-source.tex, source lines 3799--3799. Prerequisite pointer invoked at proof line(s) 3832; it names a standard fact of the same chapter that the counted argument uses, and every step of the counted proof built on it is retained verbatim in the delivered body.\par\medskip
\ifdefined\hypertarget\hypertarget{sourcepointer-2}{}\fi
\noindent\textit{Source reference.} Original lemma 0BXC: sources/stacks-chow/tags.txt, source lines 14975--14975. Prerequisite pointer invoked at proof line(s) 3846; it names a standard fact of a sibling chapter that the counted argument uses, and every step of the counted proof built on it is retained verbatim in the delivered body.\par\medskip
\ifdefined\hypertarget\hypertarget{sourcepointer-3}{}\fi
\noindent\textit{Source reference.} Original lemma 0AB7: sources/stacks-chow/tags.txt, source lines 13010--13010. Prerequisite pointer invoked at proof line(s) 3835; it names a standard fact of a sibling chapter that the counted argument uses, and every step of the counted proof built on it is retained verbatim in the delivered body.\par\medskip
\ifdefined\hypertarget\hypertarget{sourcepointer-4}{}\fi
\noindent\textit{Source reference.} Original lemma 0BFP: sources/stacks-chow/tags.txt, source lines 14392--14392. Prerequisite pointer invoked at proof line(s) 3836; it names a standard fact of a sibling chapter that the counted argument uses, and every step of the counted proof built on it is retained verbatim in the delivered body.\par\medskip
\begin{lemma}
\label{lemma-geometric-genus-normalization}
Let $k$ be a field. Let $X$ be a proper scheme of dimension $\leq 1$ over $k$.
Then
\begin{enumerate}
\item We have $g_{geom}(X/k) = g_{geom}(X_{red}/k)$.
\item If $X' \to X$ is a birational proper morphism, then
$g_{geom}(X'/k) = g_{geom}(X/k)$.
\item If $X^\nu \to X$ is the normalization morphism, then
$g_{geom}(X^\nu/k) = g_{geom}(X/k)$.
\end{enumerate}
\end{lemma}

\begin{proof}
Part (1) is immediate from Lemma \ref{lemma-bound-geometric-genus}.
If $X' \to X$ is proper birational, then it is finite and
an isomorphism over a dense open (see
Varieties, Lemmas \ref{varieties-lemma-finite-in-codim-1} and
\ref{varieties-lemma-modification-normal-iso-over-codimension-1}).
Hence $X'_{\overline{k}} \to X_{\overline{k}}$ is an isomorphism
over a dense open. Thus the irreducible components of $X'_{\overline{k}}$
and $X_{\overline{k}}$ are in bijective correspondence and the
corresponding components have isomorphic function fields.
In particular these components have isomorphic nonsingular projective models
and hence have the same geometric genera.
This proves (2).
Part (3) follows from (1) and (2) and the fact that
$X^\nu \to X_{red}$ is birational
(Morphisms, Lemma \ref{morphisms-lemma-normalization-birational}).
\end{proof}

\end{document}
